\subsection{Formulation C --- enthalpy state (v0.4)}
\label{sec:front:enthalpy}

\emph{This subsection is written at greater length than the rest of the
document. Formulation~C is the part of the model a new student is most likely to
have to modify, and it is the part where the reasoning --- not the final
formula --- is what has to be understood. Readers who only want the equations
can take Eqs.~\eqref{eq:enth-curve}, \eqref{eq:enth-ode} and
\eqref{eq:enth-exact} and skip the discussion.}

\subsubsection{The question a state variable answers}

Every transient model must decide what it \emph{remembers} from one instant to
the next. That choice --- the state variable --- is not cosmetic. It decides
which physical situations the model can represent at all, because anything the
state cannot encode is a situation the model cannot be in.

The three formulations in this section are three answers to that one question,
and they are best read as a sequence:

\begin{center}
\begin{tabular}{@{}l l l@{}}
\toprule
 & state variable & what it can represent \\
\midrule
A (v0.1) & $\delta$, from a closed form & a melt front, at one prescribed $\Tre$ \\
B (v0.2) & $\Amelt$, integrated & a melt front, driven by the actual heat flow \\
C (v0.4) & $E'$, integrated & a melt front \emph{and} a PCM temperature \\
\bottomrule
\end{tabular}
\end{center}

\noindent
Formulations A and B both store \emph{how much has melted} and nothing else.
Ask either of them ``how hot is the PCM in segment~$j$?'' and the only available
answer is $\Tm$ --- because that is all the state can say. While some solid
remains that answer is correct, and it is why the omission went unnoticed for so
long.

\subsubsection{Where Formulation B breaks, and why it is not a small error}

Consider a segment near the well head during charging. The secondary fluid
enters at $T_{4c}=\SI{160}{\celsius}$; the PCM melts at
$\Tm=\SI{150}{\celsius}$. That segment sees the hottest fluid, so it melts
first. Once its last solid is gone, the resistance network still reports a
positive heat flow --- the fluid is still \SI{10}{\kelvin} hotter than the
melting point, and nothing in Eq.~\eqref{eq:Ui} knows the segment is finished.

Formulation~B must then do one of two things, and both are wrong:

\begin{enumerate}[leftmargin=1.7em,itemsep=2pt]
\item \textbf{Keep melting.} $\Amelt$ grows past the PCM the segment actually
  contains. This is what the unguarded version did, and it produced local melt
  fractions up to $\varepsilon_{\mathrm{local}}=1.78$: the model drew latent
  heat from material that was not there.
\item \textbf{Clip and discard.} Cap $\Amelt$ at $A_{\mathrm{avail}}$ and throw
  the surplus away. Now the melt fraction is physical, but energy is not
  conserved --- and the fluid leaves the segment at the wrong temperature,
  because it gave up heat that went nowhere.
\end{enumerate}

The second is what the guarded version did, and the magnitude is the reason this
is not a rounding-level concern. At the design point the heat discarded during
\emph{discharge} reached \SI{148}{\percent} of the heat delivered. Most of what
the resistance network asked for was thrown away, and the round-trip figures
were computed on what was left.

Notice that no amount of care inside the march can fix this. The defect is in
the \emph{state}: a variable that counts melted area has no slot for ``liquid
above $\Tm$''. The physics has to be given somewhere to live.

\subsubsection{Why enthalpy, and not temperature}

The obvious repair is to carry a PCM temperature alongside the melted area. It
is also the wrong repair, and the reason is worth understanding because it
recurs throughout phase-change modelling.

At a phase change, temperature is a \emph{bad} state variable: it is not
invertible. During melting the PCM sits at $\Tm$ while its energy content
climbs by the entire latent heat, so $T$ carries no information at all
throughout the process that dominates the problem. Carrying $(T, \Amelt)$ as a
pair works, but then the code must decide at every step which of the two is
active, keep them consistent, and handle the moments when the active variable
changes.

Enthalpy is invertible. Every state of the PCM --- subcooled solid, mid-melt,
superheated liquid --- corresponds to exactly one value of $E'$, and $E'$ moves
monotonically as heat is added. So a \emph{single} scalar per segment encodes
all three regimes, and the regime is read off from the value rather than tracked
separately. Figure~\ref{fig:enthalpy-curve} is the whole idea in one picture.

\begin{figure}[htbp]
\centering
\begin{tikzpicture}[>=Stealth, scale=1.0]
  % axes
  \draw[->,thick] (0,0) -- (9.2,0) node[right] {$T_{\mathrm{pcm}}$};
  \draw[->,thick] (0,0) -- (0,5.4) node[above] {$E'$};
  % melting temperature
  \draw[dashed,gray] (4.5,0) -- (4.5,5.0);
  \node[below] at (4.5,0) {$\Tm$};
  % latent plateau markers
  \draw[dashed,gray] (0,1.6) -- (4.5,1.6) node[pos=0,left] {$0$};
  \draw[dashed,gray] (0,3.6) -- (4.5,3.6) node[pos=0,left] {$E'_{\mathrm{lat}}$};
  % the curve: solid branch, plateau, liquid branch
  \draw[very thick] (1.0,0.05) -- (4.5,1.6);
  \draw[very thick] (4.5,1.6) -- (4.5,3.6);
  \draw[very thick] (4.5,3.6) -- (8.3,5.05);
  % branch labels
  \node[align=center,font=\small] at (2.0,2.5)
    {solid,\\subcooled\\[1pt] slope $C_{s}$};
  \node[align=center,font=\small,rotate=90] at (4.05,2.6)
    {two-phase};
  \node[align=center,font=\small] at (7.1,2.9)
    {liquid,\\superheated\\[1pt] slope $C_{l}$};
  % annotate the plateau
  \draw[<->,thick,black!60] (4.9,1.6) -- (4.9,3.6)
    node[midway,right,font=\small,align=left] {latent heat\\$\rhom\hm A_{\mathrm{avail}}$};
  % annotate melt fraction
  \filldraw (4.5,2.7) circle (1.6pt);
  \draw[->,black!60] (6.2,1.5) node[right,font=\small,align=left]
    {$\Amelt = E'/(\rhom\hm)$\\here} -- (4.62,2.66);
\end{tikzpicture}
\caption{The enthalpy curve, Eq.~\eqref{eq:enth-curve}. Reading it as
$E'(T)$ it is multivalued at $\Tm$ --- which is exactly why $T$ cannot be the
state. Reading it as $T(E')$, the direction the code uses, it is single-valued
everywhere. Position on the vertical segment is the melt fraction.}
\label{fig:enthalpy-curve}
\end{figure}

\subsubsection{The state, its capacities, and their sizes}

Let $E'$ be the enthalpy per unit tube length of the PCM belonging to one
segment, measured from a datum of \emph{fully solid at $\Tm$}. The datum is
arbitrary but this choice puts the two branch boundaries at $0$ and
$E'_{\mathrm{lat}}$, which keeps the tests in the code short.

With $A_{\mathrm{avail}}=V_{\mathrm{well}}/(n_{t}L_{\mathrm{tube}})$ the PCM
cross-section a segment actually owns,
\begin{equation}
  E'_{\mathrm{lat}} = \rhom \hm A_{\mathrm{avail}},
  \qquad
  C_{s} = \rhom c_{p,s} A_{\mathrm{avail}},
  \qquad
  C_{l} = \rhom c_{p,l} A_{\mathrm{avail}},
  \label{eq:enth-caps}
\end{equation}
are the latent capacity and the two sensible capacities, all per unit tube
length. Their numerical values at the design point are worth committing to
memory, because they explain the behaviour of the whole formulation:

\begin{center}
\begin{tabular}{@{}l r l@{}}
\toprule
quantity & value & \\
\midrule
$A_{\mathrm{avail}}$      & \num{4.5411e-3} & \si{\meter\squared} \\
$E'_{\mathrm{lat}}$       & \num{2.675}     & \si{\mega\joule\per\meter} \\
$C_{s}$                   & \num{9.010e3}   & \si{\joule\per\meter\per\kelvin} \\
$C_{l}$                   & \num{1.267e4}   & \si{\joule\per\meter\per\kelvin} \\
\midrule
$C_{l}/E'_{\mathrm{lat}}$ & \num{4.74e-3}   & \si{\per\kelvin} \\
\bottomrule
\end{tabular}
\end{center}

\noindent
That last ratio is the Stefan number in disguise. One kelvin of superheat is
worth \SI{0.47}{\percent} of the latent capacity, so the sensible branches are
\emph{energetically} minor --- with the \SI{10}{\kelvin} approach available
here, at most about \SI{4.7}{\percent}. A reader could reasonably conclude that
the whole formulation is a small correction. It is not, and
\S\ref{sec:front:levelling} explains why: the sensible branches matter through
the \emph{driving temperature difference}, not through the energy they store.

The enthalpy curve and the melted area follow directly:
\begin{equation}
  T_{\mathrm{pcm}} =
  \begin{dcases}
    \Tm + E'/C_{s}, & E' < 0 \quad\text{(solid, subcooled)},\\[2pt]
    \Tm, & 0 \le E' \le E'_{\mathrm{lat}} \quad\text{(two-phase)},\\[2pt]
    \Tm + \bigl(E'-E'_{\mathrm{lat}}\bigr)/C_{l},
      & E' > E'_{\mathrm{lat}} \quad\text{(liquid, superheated)},
  \end{dcases}
  \qquad
  \Amelt = \frac{\min\!\left(\max(E',0),\,E'_{\mathrm{lat}}\right)}{\rhom \hm}.
  \label{eq:enth-curve}
\end{equation}

The clipping in $\Amelt$ is not a guard bolted on to prevent bad behaviour ---
it is the definition. Melted area saturates because a segment cannot melt PCM it
does not own, and superheat is where the extra energy goes. This is the sense in
which $\varepsilon_{\mathrm{local}}\in[0,1]$ is \emph{structural} in
Formulation~C: there is no code path that can violate it, so no test is needed.

\subsubsection{The segment equation}

The PCM in a segment exchanges heat only with the fluid passing over it (H2: no
axial conduction), so
\begin{equation}
  \frac{\mathrm{d}E'}{\mathrm{d}t}
  \;=\; q' \;=\; K\left(T_{j} - T_{\mathrm{pcm}}(E')\right),
  \qquad
  K = \frac{\md\,c_{p}\left(1-e^{-\mathrm{NTU}}\right)}{\Delta z}.
  \label{eq:enth-ode}
\end{equation}

The conductance $K$ deserves a comment, because it is where the segment march
and the melt front meet. It is not $U_{i}P_{i}$; it is the conductance
\emph{implied by the effectiveness} of Eq.~\eqref{eq:ntu}. The distinction
matters when $\mathrm{NTU}$ is large: a segment cannot extract more than the
fluid's own capacity rate $\md c_{p}$ can carry, and the factor
$\left(1-e^{-\mathrm{NTU}}\right)$ enforces that automatically. Writing $K$ this
way makes the heat taken by the PCM and the heat given up by the fluid the same
number by construction, which is what keeps energy closure structural rather
than approximate.

Two properties of Eq.~\eqref{eq:enth-ode} drive everything that follows:

\begin{itemize}[leftmargin=1.5em,itemsep=2pt]
\item It is \textbf{nonlinear}, because $T_{\mathrm{pcm}}(E')$ is piecewise.
\item It is \textbf{linear inside each branch}, because each piece of
  Eq.~\eqref{eq:enth-curve} is affine in $E'$.
\end{itemize}

\noindent
The second property is what makes an exact integration possible, and
\S\ref{sec:front:integration} takes it up. But first, why an exact integration
is not optional.

\subsubsection{Why explicit stepping fails --- and why nobody noticed until v0.4}
\label{sec:front:integration}

Look again at the three branches. On the plateau, $T_{\mathrm{pcm}}=\Tm$ is
\emph{constant}: it does not respond to $E'$ at all. The segment therefore
behaves as though it had infinite heat capacity, Eq.~\eqref{eq:enth-ode} reduces
to $\mathrm{d}E'/\mathrm{d}t = \text{constant}$, and an explicit Euler step is
not merely stable but \emph{exact}.

This is precisely the regime Formulations A and B live in permanently. Explicit
stepping was safe in v0.1 and v0.2 for a structural reason, not a lucky one ---
and that is why the time grid was never examined.

Add the sensible branches and the capacity becomes finite. Equation
\eqref{eq:enth-ode} is now a relaxation equation with time constant $C/K$, and
explicit Euler requires $\Delta t < C/K$. At the design point:

\begin{center}
\begin{tabular}{@{}l r r@{}}
\toprule
 & early ($t=\SI{1}{\second}$) & late ($t\to t_{ch}$) \\
\midrule
$\mathrm{NTU}$, segment 1        & \num{0.592}  & \num{0.083} \\
$K$ \si{[\watt\per\meter\per\kelvin]} & \num{64.3} & \num{11.5} \\
$C_{l}/K$ \si{[\second]}         & \num{197}    & \num{1106}  \\
time-grid step $\Delta t$ \si{[\second]} & \num{0.31} & \num{8491} \\
\midrule
$\Delta t\,K/C_{l}$              & \num{0.002}  & \textbf{\num{7.7}} \\
\bottomrule
\end{tabular}
\end{center}

\noindent
The grid is logarithmic, so it is fine at the start and violates the limit by
almost an order of magnitude at the end. The first implementation duly
diverged: the oscillation fed back through the fluid temperature and drove the
secondary fluid to \SI{261}{\kelvin}, some \SI{170}{\kelvin} below anything
physical.

This is a good failure to have seen once. It was loud, it was immediate, and it
was diagnosable. The dangerous version of the same error is a step ratio of
$1.5$ rather than $7.7$: still unstable in principle, but slowly enough to look
like a plausible answer.

\subsubsection{Branchwise-exact integration}

Rather than shrink $\Delta t$ --- which would multiply the cost of every run by
$\sim\!10^{3}$ --- exploit the linearity within each branch and integrate in
closed form. Holding the fluid temperature $T_{j}$ fixed across the step:

\begin{equation}
  \underbrace{E'(t+\Delta t) = E' + K\left(T_{j}-\Tm\right)\Delta t}
             _{\text{plateau: linear}},
  \qquad
  \underbrace{E'(t+\Delta t) = E'_{\mathrm{eq}}
    + \bigl(E'-E'_{\mathrm{eq}}\bigr)e^{-K\Delta t/C}}
             _{\text{sensible: exponential relaxation}},
  \label{eq:enth-exact}
\end{equation}

\noindent
with the equilibrium enthalpy and the relevant capacity
\begin{equation}
  \left(C,\;E'_{\mathrm{eq}}\right) =
  \begin{cases}
    \left(C_{s},\; C_{s}\left(T_{j}-\Tm\right)\right), & E' < 0,\\[3pt]
    \left(C_{l},\; E'_{\mathrm{lat}} + C_{l}\left(T_{j}-\Tm\right)\right),
      & E' > E'_{\mathrm{lat}}.
  \end{cases}
\end{equation}

$E'_{\mathrm{eq}}$ is the enthalpy at which the PCM would sit in equilibrium
with the fluid --- the state the segment relaxes toward and never quite reaches.
The exponential is unconditionally stable and monotone: no step size, however
large, can overshoot it. That is the property explicit Euler lacks.

\paragraph{Crossing a branch boundary.}
A step may start in one branch and end in another. The step is then split at the
crossing. Let $E'_{b}$ be the boundary being approached ($0$ or
$E'_{\mathrm{lat}}$). The crossing time is
\begin{equation}
  t^{*} =
  \begin{dcases}
    \frac{E'_{b}-E'}{K\left(T_{j}-\Tm\right)}, & \text{from the plateau},\\[8pt]
    \frac{C}{K}\,
      \ln\!\left(\frac{E'-E'_{\mathrm{eq}}}{E'_{b}-E'_{\mathrm{eq}}}\right),
      & \text{from a sensible branch},
  \end{dcases}
  \label{eq:crossing}
\end{equation}
and if $t^{*} < \Delta t$ the state is advanced to $E'_{b}$, the remaining
$\Delta t - t^{*}$ is integrated in the adjacent branch, and the test is
repeated. At most a handful of crossings can occur in one step, so the loop is
bounded.

The heat reported to the fluid is then
\begin{equation}
  q'_{j} = \frac{E'^{\,n+1}-E'^{\,n}}{\Delta t},
  \label{eq:qfromE}
\end{equation}
\emph{by definition}, not by a separate flux evaluation. Whatever the segment
gained, the fluid lost. Closure is exact to machine precision and cannot drift.

\begin{worked}[A bug worth studying: the sticky boundary]
The branch test in the code is written with closed intervals,
$0 \le E' \le E'_{\mathrm{lat}}$, so a state landing \emph{exactly} on
$E'_{\mathrm{lat}}$ is still classified as two-phase. Follow the logic: the
plateau branch computes $t^{*} = (E'_{\mathrm{lat}} - E')/\text{rate} = 0$,
finds $t^{*} < \Delta t$, advances the state to $E'_{\mathrm{lat}}$ ---
where it already is --- and subtracts zero from the remaining time. The next
iteration does the same. The segment is pinned at the boundary for ever and
never enters the liquid branch.

The symptom was not a crash. It was every segment reporting a melt fraction of
exactly $1.000$ and a superheat of exactly zero --- results that look
\emph{excellent}. The tell was a single diagnostic line printing ``sensible
energy stored: \num{0.000}'', a quantity that had no business being identically
zero.

The fix is to step just past the boundary rather than onto it,
\[
  E' \leftarrow E'_{\mathrm{lat}} + \epsilon,
  \qquad \epsilon = 10^{-9}\max\!\left(E'_{\mathrm{lat}},1\right),
\]
so the next iteration selects the liquid branch. The general lesson: when a
solver switches on the sign of a quantity, decide deliberately which side of the
boundary the equality belongs to, and make sure the state cannot come to rest
on it.
\end{worked}

\paragraph{Verification.}
The scheme was checked against a brute-force explicit integration of
Eq.~\eqref{eq:enth-ode} using \num{200000} substeps, over a charge and a
discharge, including steps that cross branch boundaries. Agreement is
$3\times10^{-13}$ relative --- machine precision. This is a genuine
verification: the two calculations share no code path beyond
Eq.~\eqref{eq:enth-curve}.

\subsubsection{Recovering the melt-layer thickness}

The resistance network needs $\delta$, not $\Amelt$. Inverting
Eq.~\eqref{eq:area-fins} --- which excludes the fin metal from the melted PCM
--- gives $\delta$ as described in \S\ref{sec:front:finmetal}, and the network
of \S\ref{sec:network} is otherwise unchanged in form.

One thing \emph{is} changed: the driving temperature. In Formulations~A and~B
the outer boundary of the resistance chain sits at $\Tm$. In Formulation~C it
sits at $T_{\mathrm{pcm}}$, which equals $\Tm$ only while some solid remains.
This single substitution is the mechanism of the next subsection.

\subsubsection{Sensible heat is self-levelling}
\label{sec:front:levelling}

Here is why a \SI{4.7}{\percent} energy term produces a first-order change in
the answer.

Consider again the well-head segment that melts first. Under Formulation~B it
finished melting and then continued to absorb heat at the full driving
difference $T_{j}-\Tm$, because its boundary temperature was pinned at $\Tm$
regardless of state. It monopolised the fluid's heat. The segments downstream,
which had not finished, received whatever was left.

Under Formulation~C the same segment starts to superheat the instant its solid
is gone. $T_{\mathrm{pcm}}$ rises, $T_{j}-T_{\mathrm{pcm}}$ collapses, its
demand falls away, and the heat continues down the well to segments that still
have solid to melt. The mechanism is a negative feedback, and it is automatic:
nothing in the code distributes the heat, and no parameter was tuned.

The effect on the distribution of melt is large:

\begin{center}
\begin{tabular}{@{}l l r@{}}
\toprule
formulation & $\varepsilon_{\mathrm{local}}$ at ten stations & spread \\
\midrule
B (latent only) & 1.78\; 0.73\; 0.69\; 0.71\; 0.77\; 0.85\; 0.94\; 1.03\; 1.12\; 1.23 & 1.321 \\
C (enthalpy)    & 1.00\; 1.00\; 1.00\; 1.00\; 1.00\; 1.00\; 1.00\; 1.00\; 1.00\; 1.00 & \textbf{0.078} \\
\bottomrule
\end{tabular}
\end{center}

\noindent
where the spread is
$\max_{j}\varepsilon_{\mathrm{local}}-\min_{j}\varepsilon_{\mathrm{local}}$ over
the \num{100} segments at end of charge. Two separate things improved. The
values above $1$ are gone, because Eq.~\eqref{eq:enth-curve} makes them
impossible. The values well below $1$ are gone too, and that is the physical
result: PCM that never melts is \emph{dead volume} --- it occupies borehole,
costs money, and carries no energy through the round trip. Formulation~C says
there is far less of it than Formulation~B implied.

\subsubsection{What Formulation C still does not do}

Three limitations, stated plainly because they bound how far the result can be
pushed.

\begin{description}[leftmargin=1.7em,itemsep=3pt]
\item[One temperature per segment.] $E'$ is a lumped quantity: all the PCM in a
  segment is at a single $T_{\mathrm{pcm}}$. There is no radial resolution
  \emph{within} a phase, so superheat appears instantaneously and uniformly
  across the whole liquid annulus rather than diffusing outward from the tube.
  The model therefore reaches its superheated state sooner than the real
  material would, and the self-levelling above is an upper bound on the effect.
\item[H1 is inherited, not repaired.] The quasi-steady melt layer of
  \S\ref{sec:front:legacy} is still assumed: the conduction profile through the
  liquid is taken as instantaneously steady. This is justified by
  $\mathrm{Ste}=\num{0.047}$ and is a good approximation here, but it is an
  approximation, and it is the same one in all three formulations.
\item[H3 is now \emph{saturated}.] Because $\varepsilon_{\mathrm{local}}\to1$
  over essentially the whole well, $\delta \to \delta_{\mathrm{merge}}$
  everywhere --- the melt fronts of neighbouring legs are touching. The annular
  conduction resistance is being evaluated at the exact limit of its validity.
  \S\ref{sec:defects:h3} takes this up, because it now carries a design
  decision.
\end{description}

Formulation~C is nonetheless the first of the three that is \emph{closed}: every
joule the resistance network delivers has a place to go, and the melt fraction
cannot leave $[0,1]$. Whatever is wrong with the answer is now wrong in the heat
transfer, not in the bookkeeping.

\emph{Source:} \code{pcm\_capacities}, \code{pcm\_state},
\code{advance\_segment}, \code{march\_h}.

